\insertImage{3a28dd01-907b-4e73-ac63-2359fb6be347.png}{Aquarelle représentant une foule de personnes, dont les visages présentent un mélange d'expressions, se tenant devant de grands bâtiments d'entreprises modernes. La foule et le bâtiment symbolisent tous deux les normes sociétales et culturelles.}

\chapter{Normes sociales et culturelles}

Vous pouvez réécrire l'organigramme de votre entreprise en une réunion. Vous ne réécrirez pas en une réunion ce que vos équipes pensent, depuis l'enfance, de ce qu'est un chef.

C'est la leçon de cette cinquième partie, et elle explique un mystère qui devrait intriguer tout le monde : pourquoi tant de libérations échouent alors que tout a été « bien fait » — la vision, les outils, l'accompagnement. Réponse : parce que l'entreprise ne se transforme pas dans le vide. Elle baigne dans une société qui a ses propres idées sur la hiérarchie, le travail et l'autorité — et ces idées ne lisent pas les livres de management.

\section{Les attentes sociétales du leadership et de la hiérarchie}

Commençons par chez nous, parce que la France est un cas d'école.

Les travaux comparatifs sur les cultures nationales classent la France parmi les pays à forte « distance hiérarchique » : on y attend du pouvoir qu'il soit centralisé, incarné, et un peu solennel\footnote{Hofstede, G. (2001). Culture's Consequences: Comparing Values, Behaviors, Institutions and Organizations Across Nations. Sage.}. Le sociologue Philippe d'Iribarne va plus loin : l'entreprise française fonctionnerait selon une « logique de l'honneur », où chacun tient à son rang, à son métier, à ses prérogatives — un héritage qui vient de bien plus loin que le capitalisme moderne\footnote{D'Iribarne, P. (1989). La logique de l'honneur: Gestion des entreprises et traditions nationales. Seuil.}. Ajoutez le culte national du diplôme et des grandes écoles, où le classement de vos vingt ans définit votre carrière, et vous obtenez une culture où l'autorité se présume plus qu'elle ne se démontre.

Maintenant, plaquez là-dessus le discours de la libération : « il n'y a plus de chefs, chacun décide ». Vous voyez le choc frontal. Ce n'est pas seulement le patron qui doit changer — ce sont les attentes de tout le monde. Y compris, et c'est le plus troublant, celles des salariés eux-mêmes : on peut détester son chef et être déstabilisé de ne plus en avoir. Les deux sentiments cohabitent très bien, car ils viennent du même endroit — une vie entière passée dans des institutions hiérarchiques, de l'école à l'armée en passant par l'administration.

Zobrist racontait d'ailleurs que le plus dur, chez FAVI, n'avait pas été de supprimer les contrôles, mais de convaincre les opérateurs que c'était vrai — qu'il n'y avait pas de piège\footnote{Zobrist, J.-F. (2012). La belle histoire de FAVI: L'entreprise qui croit que l'homme est bon. Humanisme \& Organisations.}. La confiance déclarée met des années à devenir une confiance crue. Ce délai n'est pas de la mauvaise volonté : c'est le temps que met une norme sociale à céder.

\section{La résistance culturelle au changement}

La résistance au changement a mauvaise réputation — on la traite comme une pathologie à soigner. Grave erreur de diagnostic : la résistance est souvent une réaction rationnelle de gens qui défendent quelque chose.

La sociologie des organisations l'a montré il y a longtemps : dans toute organisation, les acteurs construisent des stratégies autour des règles en place — des zones d'autonomie, des protections, des monnaies d'échange\footnote{Crozier, M., \& Friedberg, E. (1977). L'acteur et le système: Les contraintes de l'action collective. Seuil.}. Une transformation qui rebat les cartes menace ces équilibres, et ceux qui résistent ne défendent pas « le passé » par nostalgie : ils défendent leurs intérêts, leur identité professionnelle, leur maîtrise de leur propre travail. La résistance se comprend avant de se combattre — et souvent, la comprendre suffit à la dissoudre\footnote{Bouville, G. (2008). La résistance au changement: réflexion théorique et expérience vécue. Revue française de gestion, (181), 111-129.}.

Dans le cas de la libération, deux résistances culturelles spécifiques méritent le détour. La peur de la perte de contrôle, d'abord — et pas seulement chez les managers : les salariés aussi perdent un contrôle, celui de la prévisibilité. Savoir qui décide, ce qu'on attend de vous, où s'arrête votre responsabilité : c'est du confort cognitif, et le supprimer a un coût pour tout le monde.

La seconde est plus retorse : le soupçon. En France, où la relation sociale en entreprise s'est historiquement construite sur la défiance — syndicats contre patronat, contrôle contre freinage —, un patron qui annonce « je vous fais confiance, organisez-vous » déclenche une question réflexe : qu'est-ce qu'il y gagne ? Et il faut l'admettre : le soupçon n'est pas toujours infondé. Des directions ont utilisé le vocabulaire de la libération pour aplatir des organigrammes, c'est-à-dire supprimer des postes. Chaque libération cynique rend les libérations sincères plus difficiles — c'est toute la tragédie du mot.

La littérature du changement le répète depuis des décennies : sans un travail patient sur le sens, l'urgence et les premières victoires, les transformations meurent\footnote{Kotter, J. P. (1996). Leading Change. Harvard Business School Press.}. Pour une transformation qui touche à quelque chose d'aussi intime que le rapport à l'autorité, comptez double.

\section{Les valeurs sociales qui soutiennent ou entravent le modèle}

Le tableau n'est pourtant pas noir, et c'est ce qui rend le cas français si intéressant : notre culture contient aussi les ferments du modèle.

Côté soutiens : ce pays a fait une révolution au nom de l'autonomie des individus, garde une tradition philosophique qui place la liberté et la responsabilité au centre de l'existence\footnote{Sartre, J.-P. (1946). L'existentialisme est un humanisme. Éditions Nagel.}, et entretient un tissu associatif, mutualiste et coopératif dense — des générations de Français ont appris à s'organiser collectivement sans patron\footnote{Touraine, A. (1978). La voix et le regard. Seuil.}. L'aspiration à l'autonomie au travail n'est pas une importation californienne : elle a des racines locales profondes.

Côté entraves, on retrouve nos vieilles connaissances : le statut et le rang, qui font qu'une réorganisation est toujours vécue aussi comme une question d'honneur ; et un rapport anxieux à l'incertitude — la France figure parmi les cultures qui tolèrent le moins l'ambiguïté, préférant des règles, même mauvaises, à l'absence de règles\footnote{Hofstede, G. (2001). Culture's Consequences: Comparing Values, Behaviors, Institutions and Organizations Across Nations. Sage.}. Or la libération, surtout dans ses premiers mois, est une machine à produire de l'ambiguïté.

La conclusion pratique tient en une phrase : on ne libère pas une entreprise contre sa culture nationale, on la libère avec — en s'appuyant sur les valeurs qui portent et en sécurisant ce que les valeurs qui freinent ont besoin de voir sécurisé. Les libérations françaises qui réussissent ne sont pas celles qui ont importé un modèle. Ce sont celles qui ont trouvé la version française de la liberté.
